\documentclass[12pt]{article}
\usepackage[margin=1in]{geometry}
\usepackage[utf8]{inputenc}
\usepackage[english]{babel}
\usepackage{setspace}
\usepackage{url}
\usepackage{footmisc}
\usepackage{fancyhdr}

\doublespacing

\pagestyle{fancy}
\fancyhf{}
\fancyhead[R]{Vaught \thepage}  % Top right corner: "Vaught" and page number

\begin{document}

\bigskip

\section*{The Enduring Virtues of the Classical Hippocratic Oath}

For nearly 2,500 years, the \textit{Hippocratic Oath} has served as the moral compass for the medical profession\footnote{History Guild, Reference 1}. In an era when modern iterations of the oath are frequently proposed to address contemporary ethical dilemmas, it is essential to reexamine the classical version. Despite the shifting landscapes of science, economics, and politics, the original oath continues to offer unparalleled guidance in balancing intervention with restraint, defining the boundaries of healing, and reinforcing the sacred duty to serve patients. This paper argues that while modern oaths incorporate progressive viewpoints, the classical Hippocratic Oath remains superior in its holistic and time-tested approach to medical ethics.

\bigskip

\section*{Symbolism and Moral Imperative}

The oath’s invocation of deities such as Apollo, Asclepius, Hygieia, and Panacea may appear archaic at first glance\footnote{Edelstein, Reference 4}. However, these figures represent far more than mythological beings. Apollo’s embodiment of the healing arts and Asclepius’s connotation of well-being evoke a commitment to a higher, almost sacred, standard of care. While modern oaths tend to secularize these ideals by substituting religious language with generic moral imperatives, they often lose the potent imagery that reminds physicians of the gravity and nobility of their calling. The classical oath’s blend of myth and ethics serves as a timeless reminder that medicine is not merely a technical pursuit but a vocation rooted in deep moral responsibility.

\bigskip

\section*{Limits on Medical Intervention and the Pursuit of Expertise}

Critics have long pointed to certain prohibitions within the classical oath—such as the interdiction against administering a deadly drug or performing surgery without proper specialization—as outdated\footnote{SpringerLink, Reference 3}. Yet, a careful reading reveals that these clauses were designed to uphold a principle of humility and restraint. The caution against “a drug that is deadly to anyone if asked” was historically a safeguard against the misuse of medical power, rather than a definitive stance on modern practices like euthanasia. Likewise, the recommendation to defer surgery to those specifically trained in the procedure underscores a commitment to expertise and collaborative care. Modern revisions sometimes neglect these subtle lessons by focusing on isolated issues instead of the broader principle: physicians must continually recognize the limits of their own competence in the pursuit of patient well-being.

\bigskip

\section*{Addressing Modern Ethical Challenges with a Classical Foundation}

Today’s healthcare environment confronts challenges such as exorbitant costs, administrative burdens, and inequitable access to care\footnote{Financial Times, Reference 11}. Modern oaths have adapted to these issues by emphasizing patient autonomy and social justice. However, the classical oath’s enduring command to “do no harm” offers a more fundamental safeguard. Its emphasis on the sanctity of the doctor-patient relationship and the inherent risks of misusing medical authority speaks directly to the systemic challenges of modern medicine. Rather than being dismissed as antiquated, the classical oath provides a robust ethical framework that naturally extends to modern concerns. Its admonitions against exploitation and injustice implicitly demand that physicians engage not only in clinical care but also in advocacy for structural reform.

\bigskip

\section*{Conclusion}

The classical Hippocratic Oath stands as a testament to the enduring principles that continue to guide the practice of medicine. While modern iterations of the oath may incorporate contemporary issues and language, they risk diluting the profound ethical imperatives that the ancient text so powerfully embodies. In acknowledging the valuable insights of modern debates, one must ultimately recognize that the timeless wisdom of the classical oath—in its call to humility, respect for specialization, and unwavering commitment to patient welfare—remains the superior model for medical ethics today. As physicians continue to navigate an ever-changing healthcare landscape, the original oath’s legacy reminds them of their fundamental duty: to care for their patients with both skill and compassion.

\bigskip

\newpage
\begin{thebibliography}{9}

\bibitem{historyguild}
History Guild. "The Hippocratic Oath: Changing Perspectives." Accessed January 20, 2025.  
\url{https://historyguild.org/the-hippocratic-oath-changing-perspectives/}

\bibitem{einsteinmed}
Hulkower, Raphael. "The History of the Hippocratic Oath: Outdated, Inauthentic, and Yet Still Relevant." \textit{Albert Einstein College of Medicine}, Bronx, New York 10461. Accessed January 20, 2025.  
\url{https://einsteinmed.edu/uploadedFiles/EJBM/page41_page44.pdf}

\bibitem{springer}
SpringerLink. "Why the Hippocratic Oath Is Still Relevant." Accessed January 20, 2025.  
\url{https://link.springer.com/article/10.1007/s11017-011-9203-z}

\bibitem{hippocratic_oath}
Edelstein, Ludwig, translator. \textit{Hippocratic Oath.} Baltimore: Johns Hopkins University Press, 1943. Accessed via PBS website, January 20, 2025.  
\url{https://www.pbs.org/wgbh/nova/article/hippocratic-oath-today/}

\bibitem{finance}
Financial Times. "US Pharmacy Groups Marked Up Drug Prices for \$7.3bn Gain, Regulator Says." Accessed January 20, 2025.  
\url{https://www.ft.com/content/318bef8a-88fe-495a-ba4c-917fb133be6b}

\end{thebibliography}

\end{document}
